\bookchapter{Bab 17: Sambungan Telefon dengan Peranti Lain}

\chapterintro{: Jika anda menggunakan iPhone dengan Mac, AirDrop adalah fungsi pertama yang perlu anda pelajari. Pada telefon pintar, ketik Kongsi AirDrop pilih Mac anda, foto akan sampai serta-merta, tanpa menggunakan data dan tanpa Wi-Fi, menggunakan Bluetooth dan sambungan terus Wi-Fi. Terdapat juga fungsi yang mudah diabaikan dipanggil “Papan Keratan Universal” — salin teks di telefon pintar, tampal terus di Mac, seperti sihir.}

\booksection{Android + Windows (Sambung ke Windows)}

: Microsoft mempunyai App bernama “Sambung ke Windows”, terbina dalam telefon pintar Samsung, telefon pintar Android lain boleh muat turun dari kedai aplikasi. Selepas berpasangan, anda boleh terus melihat pemberitahuan telefon pintar, membalas SMS, melihat-lihat foto, malah membuat dan menerima panggilan di komputer. Saya sendiri telah menggunakannya selama dua tahun, semasa mesyuarat tidak perlu sentiasa menunduk melihat telefon pintar, pemberitahuan di skrin komputer sudah memadai.

\begin{itemize}
  \item \textbf{Kolaborasi berbilang skrin jenama}: "Kolaborasi Berbilang Skrin" Huawei dan "Xiaomi Smart Share" lebih hebat — terus menggerunkan skrin telefon pintar ke komputer, anda boleh menggunakan tetikus untuk menyeret fail dari telefon pintar ke komputer, dan juga menatal TikTok dan membalas WeChat di komputer. Salin tampal merentas peranti, seret fail, pengalaman sangat lancar.
  \item \textbf{2. Unjuk Skrin ke TV: Tonton Bersama Keluarga}Semasa pulang untuk Tahun Baru Cina, sepupu perempuan saya mahu memainkan video perjalanan graduasinya di TV. Bersusah payah untuk masa yang lama, akhirnya mendapati ia sebenarnya sangat mudah:
  \item \textbf{iPhone TV}: Tarik ke bawah Pusat Kawalan dari sudut kanan atas, ketik ikon dua kotak bertindan (Cermin Skrin), cari nama TV anda, ketik sekali dan ia akan disambungkan. Prasyaratnya ialah TV dan telefon pintar disambungkan ke Wi-Fi yang sama, dan TV menyokong AirPlay (kebanyakan TV pintar menyokongnya).
\end{itemize}

\booksection{Android TV}

: Setiap jenama mempunyai nama yang berbeza, Huawei memanggilnya “Unjuk Skrin Wayarles”, Xiaomi memanggilnya “Unjuk”, OPPO/vivo memanggilnya “Interaksi Berbilang Skrin”. Biasanya boleh ditemui dalam panel suis pantas, jika tidak jumpa, cari “unjuk skrin” dalam tetapan.

\begin{itemize}
  \item \textbf{Unjuk terus dari App video}: Jika anda hanya mahu menonton iQIYI, Bilibili, YouTube di TV, sebenarnya lebih mudah — cari ikon kecil segi empat dengan simbol Wi-Fi di antara muka main, ketik, pilih TV, video akan dipindahkan ke TV. Pada masa ini telefon pintar boleh melakukan perkara lain, video terus dimainkan di TV, tidak saling mengganggu.
  \item \textbf{Unjuk skrin tidak berjaya?}: 90\% kes adalah kerana tidak disambungkan ke Wi-Fi yang sama. 10\% lagi adalah kerana fungsi unjuk skrin TV tidak diaktifkan, pergi ke tetapan TV cari suis "Interaksi Berbilang Skrin" atau "AirPlay".
  \item \textbf{3. Dengan Jam Tangan Pintar/Gelang Pintar: Pengurus Kesihatan di Pergelangan Tangan}Saya membeli gelang pintar untuk ibu saya, pada mulanya dia berkata "Saya tidak bersukan, beli ini untuk apa". Selepas menggunakannya selama dua minggu, dia mendapati dirinya hanya tidur lebih lima jam sehari, dan mula melaraskan rutin tidurnya secara aktif. Inilah nilai paling praktikal jam tangan/gelang pintar — menukar data kesihatan yang tidak kelihatan kepada nombor yang boleh dilihat.
  \item \textbf{Penapis pemberitahuan}: Apabila telefon pintar bergetar tanpa henti dalam beg atau poket, angkat pergelangan tangan untuk melihat — sama ada WeChat dari bos atau promosi Pinduoduo, boleh menilai dalam satu saat, mesej yang tidak berbaloi untuk mengeluarkan telefon pintar biarkan sahaja.
\end{itemize}

\booksection{Penyegerakan data kesihatan}

: Jam tangan memantau kadar jantung, oksigen darah, tidur, dan langkah 24 jam, data ini akan secara automatik disimpulkan ke dalam App Kesihatan di telefon pintar. Lama-kelamaan anda boleh melihat trend anda sendiri — sama ada berjalan terlalu sedikit bulan ini, sama ada kualiti tidur minggu ini bertambah baik.

\begin{itemize}
  \item \textbf{Penyegerakan kad pengangkutan/kad akses}: Kad bas dan kad akses yang diikat pada telefon pintar, kebanyakan jenama akan menyegerakkannya secara automatik ke jam tangan. Semasa menaiki LRT tidak perlu mengeluarkan telefon pintar, pergelangan tangan dekatkan ke pintu gerbang dan terus lalu, benar-benar ada perasaan "masa depan sudah tiba".
  \item \textbf{Alat cari telefon pintar}: Jam tangan biasanya mempunyai fungsi "cari telefon pintar", ketik dan telefon pintar akan berbunyi. Saya sendiri menggunakan fungsi ini lebih kerap daripada yang dijangka — telefon pintar diletakkan sembarangan dan lupa di mana.
  \item \textbf{4. Dengan Rumah Pintar: Telefon Pintar Ialah "Suis Utama" di Rumah}pintar kedengaran sangat misteri, tetapi sebenarnya hanya tiga perkara: kawalan jarak jauh menggunakan telefon pintar, tetapkan peraturan automasi, kawalan suara.
  \item \textbf{Berikut adalah beberapa kegunaan sebenar saya sendiri, anda akan tahu betapa praktikalnya:}Penyaman udara
\end{itemize}

\booksection{: Dalam perjalanan pulang pada musim panas, hidupkan penyaman udara terlebih dahulu dengan telefon pintar, masuk ke rumah sudah sejuk. Beli soket penyaman udara tambahan lebih seratus ringgit boleh merealisasikannya, tidak perlu menukar penyaman udara.}

Lampu

: Sebelum tidur baring di atas katil, malas bangun mematikan lampu, ketik pada telefon pintar atau sebut “Hey Siri/Xiao Ai, matikan lampu”.

\begin{itemize}
  \item \textbf{Robot pembersih vakum}: Sebelum keluar rumah ketik "mula pembersihan" pada telefon pintar, pulang ke rumah lantai sudah bersih. Boleh juga tetapkan pembersihan berkala setiap hari.
  \item \textbf{Kamera}: Jika ada orang tua, kanak-kanak atau haiwan peliharaan di rumah, boleh melihat keadaan rumah bila-bila masa pada telefon pintar.
  \item \textbf{Pilih ekosistem yang mana? Ringkasan satu ayat:}Gunakan iPhone Utamakan peranti yang menyokong Apple HomeKit (lebih mahal tetapi perlindungan privasi lebih baik).
  \item \textbf{Banyak peranti dan mementingkan nisbah harga-prestasi Ekosistem Xiaomi Mi Home paling lengkap, lampu, soket, robot pembersih semuanya ada.}Gunakan telefon pintar Huawei Utamakan HarmonyOS Smart Connection, pengalaman terbaik dengan peranti Huawei.
\end{itemize}

5. Fon Kepala dan Pembesar Suara Bluetooth: Kebebasan Wayarles

\begin{itemize}
  \item Menukar dari fon kepala berwayar ke fon kepala Bluetooth, bagi saya adalah peningkatan pengalaman telefon pintar yang paling ketara. Tidak akan lagi mengeluarkan seberkas wayar berbelit dari beg, dan tidak perlu risau wayar terkait pada pemegang pintu.
  \item Perpasangan sangat mudah
  \item : Fon kepala biasanya akan memasuki mod perpasangan secara automatik kali pertama dihidupkan (lampu berkelip), telefon pintar buka Bluetooth boleh mencarinya, ketik sekali dan ia akan disambungkan. Selepas itu setiap kali dikeluarkan ia akan disambung secara automatik, tidak perlu buat apa-apa.
\end{itemize}

\booksection{Jenama sama ada bonus}

: AirPods dengan iPhone akan memaparkan pop-up tahap bateri secara automatik, dan juga boleh bertukar secara lancar antara pelbagai peranti Apple. FreeBuds Huawei dengan telefon pintar Huawei, fon kepala Xiaomi dengan telefon pintar Xiaomi juga mempunyai pop-up dan pengoptimuman serupa.

\begin{itemize}
  \item \textbf{Pembesar suara Bluetooth}: Semasa perhimpunan keluarga, berkhemah luar rumah, satu pembesar suara Bluetooth boleh menggandakan suasana. Selepas disambungkan, telefon pintar ialah DJ stand anda.
  \item \textbf{Perhatikan kependaman}: Apabila bermain permainan menggunakan fon kepala Bluetooth, mungkin ada kependaman sedikit (bunyi tembakan separuh saat baru terdengar). Jika sensitif terhadap kependaman, sama ada gunakan fon kepala berwayar, atau beli fon kepala Bluetooth yang menyokong "mod kependaman rendah".
  \item \textbf{6. Dengan Sistem Kereta: Rakan Selamat Semasa Memandu}sambil melihat telefon pintar adalah haram, dan juga boleh membawa maut. Tetapi kadang-kadang benar-benar perlu navigasi, menjawab panggilan, memainkan muzik. Pada masa ini CarPlay (Apple) dan Android Auto (Android) berguna.
  \item \textbf{Cara menyambung}: Jika kereta anda menyokongnya (kebanyakan kereta baru dalam beberapa tahun kebelakangan ini menyokong), gunakan kabel data untuk memasukkan telefon pintar ke port USB kereta, skrin kereta akan bertukar secara automatik ke antara muka CarPlay atau Android Auto. Sebahagian kereta baru menyokong sambungan wayarles, lebih menjimatkan masa.
\end{itemize}

\begin{tcbinfo}{Teknologi UWB Ultra-Wideband — "Rasa ruang" lebih tepat daripada Bluetooth}
Selain Bluetooth, telefon baharu juga ada cip UWB (Ultra-Wideband). Keupayaannya adalah\textbf{Pengukuran jarak tepat dan penentuan kedudukan ruang}— boleh kira peranti lain berapa jauh dari anda, di arah mana, ketepatan boleh capai sentimeter. Bluetooth hanya boleh beritahu “berdekatan”, UWB boleh beritahu “1.2 meter di hadapan kiri anda”. Aplikasi biasa:

\begin{itemize}
  \item \textbf{AirTag/Cari tepat}: iPhone cari AirTag papar "7 meter jauh, jalan kiri" navigasi panah, bergantung pada UWB. Bluetooth biasa hanya boleh beritahu "dalam atau luar julat".
  \item \textbf{Kunci kereta digital}: iPhone dan Android flagship dilengkapi UWB, boleh jadi kunci kereta dekat auto buka, jauh auto kunci, juga boleh bezakan kedudukan pemandu dan penumpang, anti pecah relay.
  \item \textbf{Navigasi dalaman}: Dalam pusat membeli-belah, lapangan terbang melalui beacon UWB capai navigasi "tahap lorong", lebih tepat daripada GPS.
  \item \textbf{Model yang menyokong}: iPhone 11 dan selepas semua ada cip UWB (U1/U2); Android Xiaomi, Samsung, vivo, OPPO sesetengah flagship juga telah dipasang.
\end{itemize}

\end{tcbinfo}

\begin{tcbinfo}{LE Audio dan Auracast — Standard generasi baharu audio Bluetooth}
Audio Bluetooth klasik (A2DP) yang diguna dua puluh tahun sedang diganti\textbf{LE Audio}(Audio Tenaga Rendah), ini standard audio generasi baharu selepas Bluetooth 5.2:

\begin{itemize}
  \item \textbf{Kualiti audio lebih baik: LE Audio guna codec LC3, pada kadar bit sama kualiti jauh melebihi codec SBC lama; pada kualiti sama, kadar bit lebih rendah, lebih jimat kuasa.}Penggunaan kuasa rendah: Boleh tingkat jangka hayat earphone Bluetooth kira-kira 1.5-2 kali.
  \item \textbf{Berbilang aliran: Satu telefon boleh serentak hantar audio ke berbilang earphone (senario berkongsi muzik).}Auracast Broadcast Audio: Boleh siar audio ke bilangan besar penerima tanpa had, sesuai senario awam seperti televisyen senyap gim, audio panduan lapangan terbang.
  \item \textbf{Kesan kepada pengguna: Earphone Bluetooth beli selepas 2024 semakin banyak sokong LE Audio; iOS 17+ dan Android 14+ juga sokong native. Jika beli earphone baharu, perhatikan sama ada sokong LE Audio/Auracast.}Wi-Fi 7 (802.11be) — Peningkatan kelajuan rangkaian generasi baharu\textbf{Wi-Fi 7 adalah standard Wi-Fi generasi baharu, kelajuan teori capai}Band 6GHz ditambah: Band 2.4GHz + 5GHz asal, tambah band 6GHz baharu, lebar jalur jumlah lebih besar.
  \item \textbf{Status semasa}Pengkodan 4K-QAM: Berbanding 1K-QAM Wi-Fi 6, kadar penghantaran naik 20\%.
\end{itemize}

\end{tcbinfo}

\booksection{Apa yang boleh dilakukan}

: Navigasi (jauh lebih baik daripada navigasi kereta), membuat panggilan (jawab dari stereng), mendengar muzik (App muzik dalam telefon pintar terus dimainkan di kereta), menghantar WeChat secara suara (Siri atau Xiao Ai membantu membaca mesej dan membalas).

\begin{itemize}
  \item \textbf{Kereta lama juga boleh guna}: Jika kereta anda agak lama dan tidak menyokong, beli skrin CarPlay luaran beberapa ratus ringgit, lekatkan pada papan pemuka, pengalaman hampir sama.
  \item \textbf{7. Keselamatan dan Privasi dalam Sambungan Peranti}Semakin banyak peranti disambungkan, semakin besar risiko kebocoran privasi. Beberapa perkara yang perlu diperhatikan:
  \item \textbf{Sebelum mengunjuk skrin di tempat awam, periksa dahulu sama ada terdapat mesej peribadi dalam bar pemberitahuan telefon pintar, jangan biarkan seluruh keluarga melihat kandungan yang tidak sepatutnya.}Sebelum memberikan peranti kepada orang lain atau menjualnya, nyah-ikat terlebih dahulu — pergi ke tetapan telefon pintar cari peranti tersebut, pilih "Buang" atau "Nyahpasang".
\end{itemize}

\booksection{Kamera rumah pintar, beli yang bermerek dan disulitkan, jangan beli jenama murahan puluhan ringgit. Terlalu banyak berita tentang kamera digodam dan disiarkan secara langsung.}

Mari cuba: Unjuk skrin telefon pintar ke TV, dan tetapkan satu senario pintar

\begin{itemize}
  \item Pastikan telefon pintar dan TV disambungkan ke Wi-Fi yang sama.
  \item Cari fungsi "Cermin Skrin" atau "AirPlay" di TV dan hidupkannya.
  \item Pada telefon pintar: iPhone dari Pusat Kawalan ketik "Cermin Skrin"; Android tarik ke bawah cari "Unjuk Skrin Wayarles".
\end{itemize}

\begin{tcbpractice}{Pilih satu foto kegemaran anda dan unjuk ke TV, biar ahli keluarga melihatnya.}
\begin{enumerate}
  \item Jika anda mempunyai peranti rumah pintar, cuba cipta "Mod Keluar Rumah" dalam App telefon pintar: satu ketik matikan lampu, matikan penyaman udara, mulakan robot pembersih vakum.
  \item Bab 18: Sandaran Data \& Pemindahan
  \item Seorang jiran saya, telefon pintarnya terjatuh ke dalam kolam air dan selepas dikeluarkan tidak boleh dihidupkan lagi. Yang paling menyedihkannya bukan telefon pintar itu sendiri, tetapi foto-foto di dalamnya yang tidak disandarkan selama tiga tahun — semua rekod anaknya dari lahir hingga pandai berjalan. Selepas mendengarnya hati saya sangat tidak sedap, kerana perkara ini boleh dielakkan sepenuhnya. Bab ini kita akan benar-benar memberikan insurans kepada data anda.
  \item 1. Fahami Terlebih Dahulu: Perkara Apa dalam Telefon Pintar Anda yang Hilang Akan Hilang Selamanya
  \item "Status kemandirian" data yang berbeza tidak sama, anda perlu membezakannya:
\end{enumerate}

\end{tcbpractice}

\begin{tcbinfo}{Kes 1: Pengoptimuman jangka hayat bateri praktikal}
\begin{itemize}
  \item \textbf{Kesan praktikal: Wi-Fi 7 router telah mula dijual, tetapi telefon perlu sokong. iPhone 15 Pro/16 siri, Xiaomi 14, Samsung S24 Ultra dll flagship 2024 telah sokong Wi-Fi 7; peranti lebih awal disokong melalui kemas kini OS. Jika ada router Wi-Fi 7, boleh nikmati kelajuan rangkaian lebih cepat dan lebih stabil.}: Pengguna Xiaomi 14 jangka hayat bateri rendah, separuh hari sudah tiada bateri.
  \item \textbf{Setiap telefon ada fungsi "tersembunyi" yang tidak digunakan kebanyakan orang — ia bukan Easter Egg sengaja tersembunyi oleh jurutera, tetapi ditimbus dalam lapisan tetapan dalam, atau tidak diperkenalkan dengan jelas oleh pengiklanan. Bab ini akan kumpulkan beberapa fungsi paling praktikal yang patut anda cuba.}: Tetapan -> Bateri -> Lihat ranking penggunaan kuasa -> Kenal pasti App latar belakang abnormal -> Tutup 5G tukar 4G -> Had latar belakang -> Buka mod jimat kuasa -> Cabut sarung semasa cas untuk dissipasi haba.
  \item \textbf{"Pilihan pembangun" adalah mod kawalan khas yang disediakan Android untuk jurutera, biasanya tidak ditunjukkan kepada pengguna biasa. Tetapi beberapa fungsi praktikal di dalamnya sangat berbaloi untuk dibuka.}: Jangka hayat dipanjangkan 30-50\%, data kesihatan bateri boleh dibaca.
\end{itemize}

\end{tcbinfo}
